% EgoAnchor project-specific VGTC scaffold.
% The active Chinese draft is egoanchor_cn_v1.tex. This file is a compact
% English submission skeleton kept to avoid accidentally editing the stock
% VGTC sample paper.

\documentclass[review]{vgtc}

\graphicspath{{figures/}{pictures/}{images/}{./}}

\usepackage{times}
\renewcommand*\ttdefault{txtt}
\usepackage{mathptmx}
\usepackage{booktabs}
\usepackage{amsmath}
\usepackage{amssymb}
\usepackage{graphicx}

\onlineid{0}
\vgtccategory{Research}
\vgtcinsertpkg

\title{EgoAnchor: Frame-Aligned 6DoF Object Pose Tracking with Reliability-Gated One Euro Anchor Control for Passthrough Mixed Reality}

\author{Anonymous Authors}

\teaser{
  \centering
  \fbox{\rule{0pt}{2.0in}\rule{0.95\linewidth}{0pt}}
  \caption{Planned teaser: Quest passthrough stereo capture, external 6DoF pose service, capture-time camera pose lookup by frame id, and the resulting world-consistent real-object anchor in Unity.}
  \label{fig:teaser}
}

\abstract{
EgoAnchor studies how asynchronous camera-space 6DoF object pose observations can be turned into stable, world-consistent, and recoverable real-object anchors in passthrough mixed reality. The system couples a Quest/Unity front end with an external Python perception backend and maps each returned object pose through the camera pose recorded at image capture time. The planned temporal layer is a lightweight reliability-gated One Euro anchor controller, with Kalman filtering kept as a baseline rather than as the method core: trusted observations update a One Euro filter, unreliable or stale observations are held or rejected, render-frame output uses bounded prediction, and prolonged failure triggers reacquisition. The evaluation will compare arrival-time mapping, frame-aligned raw output, low-pass, Kalman, vanilla One Euro, and the EgoAnchor modified One Euro controller using anchor-centric metrics including world-space anchor error, head-motion-induced slip, jitter, lag, latency, and recovery behavior.
}

\keywords{Mixed reality, real-object anchoring, 6DoF object pose tracking, frame alignment, One Euro filter, asynchronous sensing.}

\begin{document}
\firstsection{Introduction}
\maketitle

\section{Introduction}

MR applications consume anchor behavior, not raw pose matrices. A pose estimator may return an accurate camera-space object pose for the image it saw, but an application still needs to map that pose into the correct world frame, reject stale or unreliable observations, and maintain usable output when perception is intermittent.

EgoAnchor frames this as a pose-to-anchor problem: how should an external, asynchronous, low-rate, noisy 6DoF object pose stream become a world-consistent real-object anchor in passthrough MR? The main system claim is conservative. Frame alignment is the core mechanism; temporal control should be lightweight enough to compare directly against standard baselines.

\section{Related Work}

This section should prioritize XR/MR anchoring, registration stability, asynchronous sensing, and interactive filtering before describing 6DoF object pose estimation. One Euro filtering is cited as a speed-adaptive low-pass filter for noisy interactive input~\cite{casiez2012oneeuro}; Kalman filtering remains the canonical state-estimation baseline~\cite{kalman1960new}.

\section{Problem Definition}

For an object pose observed at image capture time \(t_i\), EgoAnchor maps the camera-space observation through the capture-time camera pose:
\[
{}^{W}\mathbf{T}_{O}(t_i) =
{}^{W}\mathbf{T}_{C}(t_i)\cdot{}^{C}\mathbf{T}_{O}(t_i).
\]
The arrival-time alternative,
\[
{}^{W}\hat{\mathbf{T}}_{O} =
{}^{W}\mathbf{T}_{C}(t_r)\cdot{}^{C}\mathbf{T}_{O}(t_i),
\]
mixes the object pose from \(t_i\) with the headset pose at return time \(t_r\), which creates visible slip under head motion.

\section{Method}

\subsection{Object-Centric Pose Stream}
Describe Quest stereo capture, camera-info remapping, prompt-guided segmentation, stereo depth, and model-based register/track/re-register. These modules are enabling components, not standalone paper contributions.

\subsection{Frame-Aligned Mapping}
Describe the \texttt{frame\_id} contract, capture-time camera pose history, OpenCV-to-Unity coordinate conversion, and world-pose composition.

\subsection{Reliability-Gated One Euro Control}
The method should stay small: score and flag gate, absolute jump rejection, position and quaternion log-space One Euro filtering, bounded render-frame prediction, static lock for residual slip, and lifecycle handling for hold, lost, and reacquire. Compare this against raw, low-pass, Kalman, and vanilla One Euro.

\section{Implementation}

Summarize the Unity/Quest front end, Python runtime, ZMQ Protobuf data plane, NATS Protobuf message/command plane, and Unity anchor runtime. Avoid legacy MessagePack or v1/v2 terminology.

\section{Evaluation}

Recommended research questions:
\begin{enumerate}
  \item Does frame-aligned mapping reduce world-space anchor error and head-motion-induced slip compared with arrival-time mapping?
  \item How do low-pass, Kalman, vanilla One Euro, and reliability-gated One Euro trade off jitter, lag, and jump suppression?
  \item How does EgoAnchor behave under occlusion, out-of-view intervals, tracking degradation, and reacquisition?
\end{enumerate}

\section{Discussion}

Discuss why camera-space pose accuracy, world-space anchor error, perceived stability, and task usability are different levels of evidence. State limitations plainly: known rigid objects, 3D model requirement, external GPU latency, and failure on reflective or poorly segmented objects.

\section{Conclusion}

The conclusion should return to the pose-to-anchor claim: passthrough MR object anchoring requires capture-time frame alignment, lightweight temporal control, lifecycle recovery, and anchor-centric evaluation.

\bibliographystyle{abbrv-doi-hyperref-narrow}
\bibliography{egoanchor_cn_refs}
\end{document}
